\answer{OA09}
对正整数 $r$，用 $\ell(r)$ 表示整除 $r$ 的最大的 $2$ 的幂，即二进制中最低非零位的值。先研究一般局面 $(r,m)$：还剩 $r$ 颗，本次最多取 $m$ 颗。下面证明
\[
 (r,m)\text{ 为必败局面}\quad\Longleftrightarrow\quad m<\ell(r).
\]
证明对剩余石子数归纳。若 $m\ge\ell(r)$，取走 $\ell(r)$ 颗。若恰好取完则获胜；否则 $r-\ell(r)$ 的最低非零位严格高于 $\ell(r)$，而对方的取数上限为 $\ell(r)$，故对方进入必败局面。

若 $m<\ell(r)$，任取合法数量 $a$，都有 $1\le a<\ell(r)$。写 $a=2^s q$，其中 $q$ 为奇数。因 $r$ 是更大的 $2$ 的幂的倍数，$r-a$ 恰含 $s$ 个因子 $2$，所以
\[
 \ell(r-a)=\ell(a)\le a.
\]
对方的上限为 $a$，故进入必胜局面。所有合法走法都把必胜局面交给对手，完成归纳。

初始局面为 $(n,n-1)$，它必败当且仅当 $\ell(n)=n$，也就是 $n$ 为$2$ 的幂。因此
\[
 \boxed{\text{先手必胜}\iff n\notin\{2^k:k\ge1\}.}
\]
必胜时第一步取 $\ell(n)$ 颗；以后每当轮到自己，就取当前剩余数量的最低非零位。上述判据保证这一步始终合法，并最终取完石子。
